\documentclass[12pt]{article}
\usepackage[T1]{fontenc}
\usepackage[utf8]{inputenc}
\usepackage{lmodern}
\usepackage{textcomp}
\usepackage{enumitem}
\usepackage{geometry}
\geometry{margin=1in}
\begin{document}
\begin{center}
Answer Key - Biology - Graduate Level Assessment 11 \
\small (Answers are ordered exactly as in the question paper.)
\end{center}

\section*{Multiple Choice}
\begin{enumerate}[label=\arabic*.]
\item \textbf{Answer:} B
\\ \textit{Options:} A) parenchyma and companion cells; B) tracheids and vessel elements; C) parenchyma and collenchyma; D) sieve tube members and companion cells; E) collenchyma and vessel elements
\\ \textit{Note:} Correct option: B - tracheids and vessel elements
\item \textbf{Answer:} B
\\ \textit{Options:} A) increases as H+ concentration increases; B) decreases as blood pH decreases; C) decreases in resting muscle tissue; D) increases as blood pH decreases; E) increases as O 2 concentration decreases
\\ \textit{Note:} Correct option: B - decreases as blood pH decreases
\item \textbf{Answer:} A
\\ \textit{Options:} A) The temperature during embryonic development determines the sex.; B) The presence of an X or Y chromosome in the mother determines the sex of the offspring.; C) The presence of the X chromosome determines the sex.; D) The sex is determined by environmental factors after birth.; E) The presence of two Y chromosomes determines that an individual will be male.
\\ \textit{Note:} Correct option: A - The temperature during embryonic development determines the sex.
\item \textbf{Answer:} A
\\ \textit{Options:} A) They have flowers that bloom for several years before producing seeds; B) They have a unique leaf structure; C) They can survive in extremely cold climates; D) They produce large seeds; E) Their seeds can remain dormant for up to a century before germinating
\\ \textit{Note:} Correct option: A - They have flowers that bloom for several years before producing seeds
\item \textbf{Answer:} A
\\ \textit{Options:} A) produce ATP; B) create carbon dioxide; C) make FADH 2; D) synthesize proteins; E) make glucose
\\ \textit{Note:} Correct option: A - produce ATP
\end{enumerate}

\section*{True / False}
\begin{enumerate}[label=\arabic*.]
\setcounter{enumi}{5}
\item \textbf{Answer:} False
\\ \textit{Options:} A) True; B) False
\\ \textit{Note:} Laparoscopic antireflux surgery is a safe and effective treatment for GERD even in elderly patients, warranting low morbidity and mortality rates and a significant improvement of symptoms comparable to younger patients.
\item \textbf{Answer:} True
\\ \textit{Options:} A) True; B) False
\\ \textit{Note:} A greater association than hitherto acknowledged, between ascitis volume and anthropometric measurements from one side, and long-term rehospitalization and mortality from the other, was demonstrated in male stable alcoholic cirrhotics. Further studies with alcoholic and other modalities of cirrhosis including women are recommended.
\item \textbf{Answer:} False
\\ \textit{Options:} A) True; B) False
\\ \textit{Note:} The rabbit is a good model to be used in training of surgery, with a low morbi-mortality, able to be anesthetized intramuscularly, with no need of pre-operative fasting and does not present hypoglycemia even with the extended fasting period.
\item \textbf{Answer:} True
\\ \textit{Options:} A) True; B) False
\\ \textit{Note:} Postoperative numbness occurs in most patients receiving nasal microfat injections. Partial to complete recovery of nasal tip sensation can be expected to occur over a 3-month period.
\item \textbf{Answer:} True
\\ \textit{Options:} A) True; B) False
\\ \textit{Note:} A fetal anatomic survey on follow-up sonograms may identify unanticipated fetal anomalies, especially when the indication is for fetal growth.
\end{enumerate}

\section*{Long Form Response}
\begin{enumerate}[label=\arabic*.]
\setcounter{enumi}{10}
\item \textbf{Answer:} Procaryotic and Eucaryotic. Explain why this is the correct answer and provide supporting evidence. Reference key facts or concepts that support this choice.
\\ \textit{Note:} Long-form derived from MCQ; award credit for stating the correct idea and giving supporting reasoning.
\item \textbf{Answer:} Living things are characterized by a specific, complex organization, the ability to consume energy, movement, response to physical or chemical changes, growth, reproduction, and adaptability.. Explain why this is the correct answer and provide supporting evidence. Reference key facts or concepts that support this choice.
\\ \textit{Note:} Long-form derived from MCQ; award credit for stating the correct idea and giving supporting reasoning.
\end{enumerate}

\vfill
\noindent\textit{End of Answer Key}
\end{document}